% =============================================================================
% New \section{} for NOETHER_paper.tex
%
% Insertion point:
%   AFTER current \section{Discussion and threats to validity} starts at
%   line 994; BEFORE that, between current \S6 (Cross-domain demonstration,
%   line 553) and current \S7 (Discussion, line 994).
%
% Renumbering: current \S7 -> \S8 (Discussion); current \S8 -> \S9
% (Conclusion). Internal \ref{} are unaffected because they go through
% \label{}.
%
% Scope: programs with explicit mathematical-physics equations
% (NOETHER's pre-declared in-scope class, see \S\ref{subsec:scope}).
% Comparator: GenMorph GP-evolved Set G (Ayerdi et al. 2023), the SOTA
% automated MR identification pipeline on Java SUTs.
%
% Two arms: (i) the GenMorph 23-method utility-method baseline (see
% supplementary S5; mixed in/out-of-scope substrate; reported as a
% reference); (ii) the algebra-rich subset of the GenMorph benchmark
% selected by a pre-registered structural-coverage criterion (this is
% the in-scope evaluation; see supplementary S7).
%
% Future-work arms: Set M (MR-Scout), Set L (LLM-prompted), Set B
% (literature) on the same substrate. Honest scope: this section
% delivers the Set N vs Set G arm and explicitly defers the others.
%
% Foregrounded claims (per docs/paper_claims_summary.md):
%   B  STRONG  L*-block blindness empirically confirms the operative
%              role of the eight-block decomposition.
%   C  STRONG  Set N is defined and runnable on two SUTs (instance
%              method, boolean predicate) where Set G is structurally
%              absent, evidencing coverage extension.
%   D  STRONG  Scope declaration is empirically discriminating
%              (utility-method baseline -> algebra-rich gain).
%   F  STRONG  GP independently rediscovers Set N's algebraic core on
%              \texttt{midpoint}.
%   E  REFRAMED  Reframed from "Set N > Set G" to
%              "competitive parity with coverage extension".
% =============================================================================

\section{Comparative empirical evaluation against SOTA on programs with explicit mathematical-physics equations}
\label{sec:empirical-vs-sota}

The previous two sections establish that NOETHER \emph{generates} MRs
in two domains with structurally distinct operator algebras: Boltzmann
transport (\S\ref{sec:reactor}) and equivariant machine learning
(\S\ref{sec:cross-domain}). Generation alone does not establish that
the generated MRs are competitive against an automated SOTA pipeline
on the substrate that pipeline was designed for. This section reports
that comparison on the in-scope class declared at
\S\ref{subsec:algebra-induced-MRs}, namely programs with explicit
mathematical-physics equations.

The comparator is GenMorph~\cite{Ayerdi2023GenMorph}, the
genetic-programming MR identification pipeline that targets Java SUTs
and reports kill-rate gains against PIT mutants. We evaluate on two
arms: a utility-method reference arm using GenMorph's published
23-method benchmark (mixed in-scope and boundary substrate; supplementary
S5), and a pre-registered algebra-rich arm restricted to subjects whose
induced operator algebra has at least one non-empty block beyond $G$
(supplementary S7). The Set M (MR-Scout-mined~\cite{MRScout2024}),
Set L (LLM-prompted), and Set B (literature) arms named at
\S\ref{para:comp-eval-protocol} remain as committed future work.

\subsection{Scope, claims, and pre-registration}
\label{subsec:empirical-scope}

NOETHER's scope declaration (\S\ref{subsec:algebra-induced-MRs} and
\S\ref{subsec:principal-limitation}) is that the framework derives
useful MRs for programs whose induced algebra $\mathcal{A}_P$ has a
non-trivial decomposition under the eight-block hypothesis. The
declaration is a \emph{falsifiable} prediction: if the framework's
kill-rate gain over a non-algebra-derived baseline does not
discriminate between in-scope and out-of-scope substrates, the
operative-mechanism reading collapses into rhetoric. This section
tests that prediction directly.

Concretely, the section evaluates four claims, each tied to a single
section of the framework's theory:
\begin{description}[leftmargin=*,nosep]
  \item[Claim B (mechanism operativeness).] The eight-block
        decomposition (\S\ref{subsec:decomposition}) is empirically
        operative: per-block kill-rate patterns match per-block
        algebraic predictions, including a falsifiable prediction
        about $\mathcal{L}^{*}$-block blindness on
        homogeneity-preserving mutators.
  \item[Claim C (coverage extension).] Algebra-derived MRs are
        defined on programs (instance methods, boolean classifiers)
        on which the SOTA GP pipeline is structurally absent;
        coverage extension is independent of kill-rate statistics.
  \item[Claim D (scope discriminance).] The kill-rate gain over the
        non-algebra-derived baseline is larger on the in-scope
        algebra-rich substrate than on the wider utility-method
        substrate, confirming that the scope declaration is
        empirically discriminating.
  \item[Claim F (cross-pipeline rediscovery).] GenMorph's GP
        independently rediscovers Set~N's algebraic core on at
        least one SUT, corroborating that algebraic structure is the
        operative generator of effective MRs.
\end{description}

We do \emph{not} claim Set~N $>$ Set~G in pooled head-to-head
kill-rate. The data does not support that claim: pooled $\Delta$ is
$-0.071$ absolute against Set~N at $n = 70$, McNemar exact
$p = 0.359$, and $n$ is underpowered for an $\alpha = 0.05$ verdict.
The supported reading is \emph{competitive parity with coverage
extension}, made precise in \S\ref{subsec:reframing}.

\paragraph{Pre-registration of the algebra-rich SUT criterion.}
The selection rule for the in-scope arm is committed in
\texttt{configs/d4j\_algebra\_rich\_criterion.json} on branch
\texttt{feat/d4j-algebra-rich} of the experiment repository
\emph{before} any new evaluation data existed. The criterion
references package roots (\texttt{org.apache.commons.math3.ode},
\texttt{*.linear}, \texttt{*.transform},
\texttt{*.analysis.solvers}, \texttt{*.distribution},
\texttt{*.optim}, \texttt{*.fitting}, \texttt{*.stat.regression},
\texttt{*.complex}, \texttt{*.geometry}, \texttt{*.fraction} and their
\texttt{math.*} legacy counterparts), per-package block-coverage
hypotheses, and method-signature constraints; it contains no bug-id
and no kill-rate references and cannot be tuned by evaluation
outcomes. The git timestamp chain (criterion $\rightarrow$ inscope
filter $\rightarrow$ Set~N derivation $\rightarrow$ Set~G GP rerun
$\rightarrow$ pooled M1) is the auditable proof of pre-registration.

\subsection{Substrate, MR sets, and metrics}
\label{subsec:empirical-setup}

\paragraph{Substrate.}
PIT 1.7.4~\cite{Coles2016PIT} with the default mutator configuration
serves as the shared substrate. Each (subject, MR) pair is evaluated
through (i) JUnit test-class codegen, (ii) a 2-pass surefire
green-suite filter that drops MR instances flaky on the unmutated
SUT, and (iii) PIT mutation testing parsed to a per-mutant kill
vector. The pooled M1 metric is the union-kill rate across MRs in the
set, reported with Wilson 95\% confidence intervals; paired
comparisons use the McNemar exact two-sided test.

\paragraph{MR sets.}
\textit{Set~N (NOETHER-derived)}: hand-derived from the operator
algebra of each SUT through CONSTRUCT-MP, expressed in GenMorph's
JIR/JOR DSL. The derivation uses only the SUT's signature and its
algebraic structure; it does not consult mutation outcomes, test
inputs, or PIT operator definitions.

\textit{Set~G (GenMorph GP-evolved)}: harvested by rerunning
GenMorph's \texttt{genmorph.py gen} mode at seed $=11$. The 1\,min
GAssert budget used for parallel scheduling handicaps Set~G against
the upstream-default 30\,min; the budget asymmetry is conservative in
the direction of \emph{against} any parity claim made in Set~N's
favour.

\paragraph{Reference and main arms.}
The reference arm uses GenMorph's published 23-method
benchmark~\cite{Ayerdi2023GenMorph}, a mixed substrate of utility
methods (\texttt{gcd}, \texttt{sin}, \texttt{capitalize}, \ldots).
The main arm uses the 10 SUTs hand-derived from the pre-registered
algebra-rich criterion, inlined into a single \texttt{MathSignalClass}
or, for the instance-method case, \texttt{ComplexSignal}, so that PIT
mutates the algorithm directly rather than a delegator.

\subsection{Reference arm: utility-method baseline}
\label{subsec:utility-arm}

The utility-method arm establishes the baseline against which the
algebra-rich arm is read. Pooled across the 23 GenMorph subjects
($\approx 575$ PIT mutants), Set~N's M1 is $0.288$ and Set~G's M1 is
$0.363$; McNemar exact $p = 0.176$ at the 30\,min GAssert
budget~\cite{Ayerdi2023GenMorph}. The full per-subject breakdown is
reported in supplementary S5; here we use only the pooled figure.

\textit{Reading.} On a substrate that is mixed in-scope (\texttt{gcd},
\texttt{sin}) and out-of-scope (\texttt{capitalize}, parser-style
predicates) for NOETHER, Set~G is directionally ahead of Set~N by
$0.075$ absolute, the difference is not significant at $\alpha=0.05$,
and Wilson 95\% CIs overlap. The reading we attach to this arm is
deliberately weak: it is the boundary line below which a
scope-conformant evaluation \emph{must} improve, otherwise the scope
declaration is rhetorical.

\subsection{Main arm: pre-registered algebra-rich evaluation}
\label{subsec:algebra-rich-arm}

\paragraph{Subjects and MR yield.}
Ten SUTs span the symmetry, translation/period, scaling, and
identity blocks of $\mathcal{A}_P$: \texttt{midpoint},
\texttt{exactLog2}, \texttt{isSequence} (boolean predicate),
\texttt{clamp}, \texttt{signum}, \texttt{ComplexSignal.add} (instance
method), \texttt{gcdSig}, \texttt{lcmSig}, \texttt{hypotSig},
\texttt{powerSig}. Set~N contains 30 hand-derived NOETHER MRs, two to
four per SUT, one or more per non-empty block of each SUT's induced
algebra. The GP rerun yields 26 Set~G MRs across 8 SUTs at the 1\,min
GAssert budget; two SUTs return Set~G structurally not available, as
detailed next.

\paragraph{Set G is structurally absent on two of the ten SUTs.}
\label{para:set-g-na}
GenMorph's GP pipeline fails on two SUTs for reasons internal to its
plumbing rather than the underlying mutation-testing problem:
\begin{itemize}[nosep,leftmargin=*]
  \item \texttt{ComplexSignal.add} (instance method) triggers an
        XStream \texttt{NullPointerException} in
        \texttt{MethodTestTransformerConfig} when deserialising the
        receiver object;
        \texttt{ch.usi.gassert.data.types.AbstractCollectionConverter}
        in the upstream snapshot does not handle user-defined
        receivers.
  \item \texttt{isSequence} (boolean predicate) yields GP-evolved JOR
        strings containing dangling \texttt{<} operators, e.g.\
        \texttt{((\ldots <))}, that do not parse as Java; the GP
        grammar for predicate-output SUTs is degenerate at this scale.
\end{itemize}
Both SUTs admit Set~N MRs (commutativity for
\texttt{ComplexSignal.add}, monotonic-window translation for
\texttt{isSequence}) and run to completion. Together the two SUTs
cover 8 of the 70 PIT mutants under evaluation, on which only Set~N
is defined. This is structural-coverage extension reported
independent of any kill-rate statistic and is the evidential basis of
claim~C.

\paragraph{Pooled head-to-head result.}
\begin{table}[h]
\centering
\caption{Per-SUT PIT mutation kill counts on the 10 algebra-rich Java
SUTs of \S\ref{subsec:algebra-rich-arm}. Set~G is structurally absent
on \texttt{ComplexSignal.add} (XStream receiver-deserialisation
failure) and \texttt{isSequence} (degenerate GP grammar on boolean
predicates); see \S\ref{para:set-g-na}. Pooled across the 8
head-to-head SUTs ($n = 70$ mutants): Set~N $= 34$, Set~G $= 39$,
McNemar exact $p = 0.359$ at the 1\,min GAssert budget. Wilson 95\%
confidence intervals on the pooled rates overlap.}
\label{tab:algebra-rich-pooled}
\small
\begin{tabular}{lccccc}
\toprule
SUT & mutants & Set N kills & Set G kills & $\Delta$ & note \\
\midrule
\texttt{ComplexSignal.add} &  3 & 2 & N/A &      & instance method \\
\texttt{midpoint}          &  3 & 3 & 3   & $0$  & tie \\
\texttt{exactLog2}         & 10 & 4 & 0   & $+4$ & N narrow \\
\texttt{isSequence}        &  5 & 0 & N/A &      & boolean predicate \\
\texttt{clamp}             &  7 & 3 & 6   & $-3$ & G \\
\texttt{signum}            &  6 & 4 & 4   & $0$  & tie \\
\texttt{gcdSig}            &  9 & 6 & 5   & $+1$ & N narrow \\
\texttt{hypotSig}          &  4 & 2 & 4   & $-2$ & G \\
\texttt{lcmSig}            & 11 & 4 & 7   & $-3$ & G \\
\texttt{powerSig}          & 12 & 8 & 7   & $+1$ & N narrow \\
\midrule
\textbf{pooled (n=70)}     &    & \textbf{34} & \textbf{39} & $-5$ & ns \\
\bottomrule
\end{tabular}
\end{table}

Across the eight head-to-head SUTs, Set~N kills 34 PIT mutants of 70
and Set~G kills 39 of 70 (Table~\ref{tab:algebra-rich-pooled}). The
pooled M1 kill rates are $0.486$ for Set~N (Wilson 95\% CI
$[0.372,\,0.600]$) and $0.557$ for Set~G ($[0.441,\,0.668]$); McNemar
exact gives $p = 0.359$. The 95\% CIs overlap and the test fails to
reject the null at $\alpha = 0.05$. \textbf{$n = 70$ across $10$ SUTs
is underpowered for a paired hypothesis test at $\alpha = 0.05$}; the
absence of a significant difference is reported as descriptive
evidence, not as confirmation of equivalence. The directional gap
($-0.071$ absolute against Set~N) is in the same direction as the
utility-method arm but smaller in absolute size and substantially
smaller in relative size.

\paragraph{Claim D: the scope declaration is empirically discriminating.}
The pre-declared scope predicts that Set~N's M1 should improve on
in-scope substrates relative to the mixed utility-method baseline.
The data shows Set~N moves from $0.288$ (utility-method baseline) to
$0.486$ (algebra-rich arm), a $+69\%$ relative gain. Set~G also
improves, from $0.363$ to $0.557$, a $+53\%$ relative gain. Both
methods benefit from the algebraically richer substrate, as expected
of any MR-based testing on programs with explicit mathematical-physics
equations; Set~N benefits more in relative terms, closing the
absolute gap from $-0.075$ to $-0.071$ and the relative gap from
$-21\%$ to $-13\%$. This is the scope-discriminance prediction, and
it is consistent with the data.

\subsection{Per-block kill pattern: the eight-block decomposition is empirically operative}
\label{subsec:per-block-kill}

Set~N MRs are tagged by the NOETHER block they instantiate, so the
70-mutant outcome decomposes into per-block per-MR kill rates rather
than a single aggregate. The pattern matches the algebraic-mechanism
prediction in three respects, including one prediction that is sharp,
falsifiable, and confirmed.

\textit{$T^{*}$ block (translation, period).} Translation MRs carry
the highest single-MR kill rates on the numeric SUTs:
\texttt{midpoint} \texttt{T\_shift} 3/3, \texttt{powerSig}
\texttt{T\_exp\_step} 8/12, \texttt{exactLog2} \texttt{T\_double}
4/10, \texttt{clamp} \texttt{T\_shift} 3/7. Translation invariance is
the property most sensitive to PIT's arithmetic-operator swaps, and
$T^{*}$-derived MRs detect this fingerprint reliably.

\textit{$G$ block (group, symmetry).} Symmetry MRs are moderately to
highly effective when the SUT exposes the appropriate symmetry:
\texttt{signum} \texttt{G\_negate} 4/6, \texttt{gcdSig}
\texttt{G\_swap} 5/9, \texttt{midpoint} \texttt{G\_swap} 1/3,
\texttt{ComplexSignal.add} \texttt{G\_swap} 2/3.

\textit{$\mathcal{L}^{*}$ block (linearity, scaling): a uniformly
near-zero kill rate, theoretically predicted.} The scaling MRs of the
form $f(\lambda \mathbf{x}) = \lambda \cdot f(\mathbf{x})$ kill 0 of
3, 0 of 7, 0 of 9, 0 of 11, and 2 of 4 mutants on \texttt{midpoint},
\texttt{clamp}, \texttt{gcdSig}, \texttt{lcmSig}, and
\texttt{hypotSig} respectively. The flat near-zero pattern is the
algebraic-mechanism claim's sharpest empirical witness: PIT's default
mutator set ($+/-/\times/\div$ swaps, \texttt{return zero},
\texttt{return one}, conditional negation) preserves homogeneity of
degree~$1$, so a paired MR of the form
$f(\lambda \mathbf{x}) = \lambda \cdot f(\mathbf{x})$ cannot in
principle distinguish the original from such mutants. The
$\mathcal{L}^{*}$-block MRs are blind to PIT's homogeneity-preserving
mutants \emph{by construction}, and the data confirms this in 5 of 6
SUTs where $L_{\mathrm{scale}}$ applies. We name this finding
\emph{$\mathcal{L}^{*}$-block blindness} and treat it as direct
empirical witness for the conservation-law-to-MR-blindness
correspondence articulated at \S\ref{subsec:algebra-induced-MRs}: a
block whose generator is a continuous symmetry of the mutator set
contributes MRs that are necessarily silent on mutator-induced
defects of that block.

\textit{$\mathcal{I}^{*}$ block (idempotence, identity).} Idempotence
MRs kill few mutants under the paired-MR DSL: \texttt{I\_idem}
$\approx 0$ across most SUTs; \texttt{powerSig} \texttt{I\_zero\_exp}
2/12 (a degenerate single-input case). This reflects an expressivity
limit of the JIR/JOR shape, not an absence of $\mathcal{I}^{*}$
structure in the SUTs.

\textit{Reading.} The block-by-block predictions (high $T^{*}$
sensitivity on numeric SUTs, moderate $G$ correlated with SUT
symmetry, near-zero $\mathcal{L}^{*}$ on
homogeneity-preserving mutators, low $\mathcal{I}^{*}$ from the DSL)
match the empirically observed per-block per-MR rates. The
$\mathcal{L}^{*}$-block prediction is the sharpest because it is
\emph{quantitative} (a near-zero rate, not a directional preference)
and \emph{ex-ante} derivable from PIT's mutator semantics composed
with the block's algebraic generator. Confirmation of this prediction
is what we read as direct evidence for claim~B.

\subsection{Cross-pipeline rediscovery: GP recovers Set N's algebraic core}
\label{subsec:rediscovery}

On \texttt{midpoint}, GenMorph's GP independently evolves three MRs
that map onto Set~N's blocks (Table~\ref{tab:rediscovery}).

\begin{table}[h]
\centering
\caption{Cross-pipeline rediscovery on \texttt{midpoint}. GP-evolved
MRs from the Set~G arm versus their algebra-derived Set~N
counterparts.}
\label{tab:rediscovery}
\small
\begin{tabular}{lll}
\toprule
GP-evolved MR (Set G) & Set N counterpart & Block \\
\midrule
\texttt{SwitchParams??1@2}              & \texttt{G\_swap}  (commutativity)        & $G$ \\
\texttt{NumericAddition?1.000000?1}     & \texttt{T\_shift} (translation by $1$)   & $T^{*}$ \\
\texttt{NumericMultiplication?0.500000?2} & \texttt{L\_scale} (scaling, approx.)   & $\mathcal{L}^{*}$ \\
\bottomrule
\end{tabular}
\end{table}

Set~N is derived a-priori from the operator algebra of
\texttt{midpoint}; Set~G is searched by mutation-killing fitness with
no algebraic structure as input. The two pipelines converge on the
same algebraic primitives. We read this as direct corroboration of
the framework's central claim that algebraic structure is the
operative generator of effective MRs, not an analyst's after-the-fact
rationalisation. This is the evidential basis of claim~F.

\subsection{Reading: competitive parity with coverage extension}
\label{subsec:reframing}

The pooled head-to-head comparison does not support a
Set~N $>$ Set~G claim: Set~G is directionally ahead by $0.071$
absolute, McNemar fails to reject the null, and $n = 70$ is
underpowered. What the data \emph{does} support is a two-part claim:

\begin{enumerate}[nosep,leftmargin=*]
  \item \emph{Competitive parity on the head-to-head substrate.} On
        the eight SUTs both methods can address, the pooled kill
        rates are within McNemar non-significance at the 1\,min
        Set~G budget. Per-SUT, Set~N narrowly wins 4 SUTs
        (\texttt{midpoint}, \texttt{exactLog2}, \texttt{gcdSig},
        \texttt{powerSig}), Set~G clearly wins 3 SUTs
        (\texttt{clamp}, \texttt{hypotSig}, \texttt{lcmSig}), 1 SUT
        ties non-trivially (\texttt{signum}), and 1 SUT ties at zero
        (\texttt{isSequence}, where the paired-MR DSL is
        structurally inadequate for boolean classifiers).
  \item \emph{Structural coverage extension.} Set~N is defined and
        runnable on the two SUTs (\texttt{ComplexSignal.add},
        \texttt{isSequence}) where GenMorph's GP pipeline is
        structurally absent. The two SUTs together cover $11.4\%$
        of the mutant denominator under evaluation, and on these
        SUTs the algebra-derived method is currently the only
        option among the compared baselines.
\end{enumerate}

We frame the contribution of \S\ref{subsec:algebra-rich-arm} as
\emph{competitive parity with coverage extension}, not as a
head-to-head superiority claim. The structural-extension component is
non-statistical and will not flip with sample size; the parity
component is reported under explicit budget asymmetry (1\,min Set~G
vs 30\,min upstream default) and explicit underpowering, both of
which are conservative in the direction of \emph{against} the parity
reading.

\subsection{Threats to validity and committed future work}
\label{subsec:empirical-threats}

\paragraph{(a) Set G budget asymmetry.}
GenMorph's published configuration is 30\,min GAssert; we ran at
1\,min for parallel scheduling. The 1\,min budget handicaps Set~G in
the conservative direction relative to any parity claim made in
Set~N's favour. A 30\,min rerun is option (a) of the future-work
table below; expected effect is a widening of Set~G's lead on the
head-to-head SUTs, which would weaken the parity reading but would
not invalidate the structural-extension finding.

\paragraph{(b) Sample size.}
$n = 70$ across 10 SUTs is underpowered for a head-to-head paired
verdict at $\alpha = 0.05$. The contribution is reframed as
competitive parity rather than superiority, with the parity reading
qualified as descriptive evidence. Extending to the full 38 in-scope
D4J subjects identified by the pre-registered criterion remains future
work; current evidence does not allow a prediction about the McNemar
verdict at $n > 200$.

\paragraph{(c) Set G structural absence is upstream-snapshot-relative.}
The two N/A SUTs reflect the state of GenMorph upstream at our
snapshot. An upstream patch (an XStream Converter for user-defined
receivers and a richer JOR grammar for boolean-output SUTs) would
restore Set~G coverage on these SUTs. The structural-extension claim
is reported with that future-state qualifier; the claim does not
assert that algebra-derived MRs are the only kind \emph{that could in
principle} be defined on these SUTs, only that they are the only kind
a current automated pipeline \emph{does} produce.

\paragraph{(d) Substrate selection.}
The evaluation restricts itself to programs with explicit
mathematical-physics equations, by design. The kill-rate gain over
the utility-method reference is expected by the framework's scope
declaration (\S\ref{subsec:algebra-induced-MRs}); it is reported as
scope-discriminance evidence (claim~D), not as a generalisable
performance improvement on every Java SUT. The reference arm itself
makes that scope explicit: on the wider mixed substrate, Set~N is
behind Set~G; on the in-scope substrate, Set~N achieves parity with
coverage extension. The paper does not generalise beyond this scope.

\paragraph{(e) Single SOTA comparator.}
The protocol of \S\ref{para:comp-eval-protocol} commits Set~N
against four baselines (Set~M, Set~G, Set~L, Set~B). This section
delivers the Set~G arm; Set~M (MR-Scout-mined~\cite{MRScout2024}),
Set~L (LLM-prompted), and Set~B (literature) remain as committed
future work, called out under (b)/(d) of Table~\ref{tab:future-work}.
The single-comparator scope is the primary reason this section's
evidence is positioned as \emph{competitive parity} rather than as a
broader empirical superiority claim.

\begin{table}[h]
\centering
\caption{Committed future work for \S\ref{sec:empirical-vs-sota}.}
\label{tab:future-work}
\small
\begin{tabular}{lll}
\toprule
Direction & Cost (estimated) & Expected effect on the reading \\
\midrule
(a) GP rerun at 30\,min GAssert budget                       & $\approx 30$\,min wall (parallel) & Probably widens Set~G's lead; conservative against parity \\
(b) Extend to all 38 in-scope D4J subjects                    & $\approx 10$\,h human + $\approx 30$\,min compute & $n > 200$; trend uncertain at this sample size \\
(c) Patch GenMorph upstream (instance methods, boolean JORs)  & 1--2 days                         & Restores Set~G on the two N/A SUTs; weakens structural framing of claim~C \\
(d) Add Set~M (MR-Scout) and Set~L (LLM) arms                  & $\approx 1$ day per arm           & Completes the comparator protocol of \S\ref{para:comp-eval-protocol} \\
\bottomrule
\end{tabular}
\end{table}

\subsection{Summary of evidence}
\label{subsec:empirical-summary}

The section delivers four lines of evidence:

\begin{itemize}[nosep,leftmargin=*]
  \item Claim~B (mechanism operativeness): supported strongly by the
        per-block kill pattern of \S\ref{subsec:per-block-kill},
        with $\mathcal{L}^{*}$-block blindness as the sharpest
        ex-ante prediction confirmed.
  \item Claim~C (coverage extension): supported strongly by the two
        Set~G N/A SUTs of \S\ref{para:set-g-na}, on which only
        Set~N is defined.
  \item Claim~D (scope discriminance): supported by the
        $0.288 \rightarrow 0.486$ Set~N gain and the
        $0.363 \rightarrow 0.557$ Set~G gain when restricted to the
        in-scope algebra-rich substrate.
  \item Claim~F (cross-pipeline rediscovery): supported by GP's
        independent rediscovery of Set~N's algebraic core on
        \texttt{midpoint} (\S\ref{subsec:rediscovery}).
\end{itemize}

The section does \emph{not} deliver a head-to-head superiority claim
(claim~E in the experiment audit at supplementary~\ref{S7-d4j}); the
reading delivered instead is competitive parity with coverage
extension, qualified by explicit underpowering and explicit budget
asymmetry both running in the conservative direction against parity.
The full per-SUT outcome matrix, the per-block kill table for all 30
Set~N MRs, the GP raw output and the \texttt{.jir}/\texttt{.jor}
translation, the McNemar $b/c$ counts on the eight head-to-head SUTs,
the reproducibility checklist, and the test-gate harness
(\texttt{tests/run.sh}) are in supplementary~\ref{S7-d4j}.
